Mixed-integer programming (MIP) remains one of the most expressive paradigms for decision optimization because it unifies continuous resource allocation and discrete strategic choice in a single mathematical model \cite{bertsimas2005optimization,morrison2016branch}. This modeling power explains why MIP appears in production planning, scheduling, routing, energy systems, finance, network design, and increasingly in data-driven decision pipelines. The same integrality structure that makes MIP valuable, however, also creates its central computational difficulty. Feasible regions become fragmented, strong linear-relaxation solutions can remain stubbornly fractional, and promising trajectories may fail at the last step from near-feasible to strictly feasible solutions. In many operational settings, that final transition matters more than early claims about ultimate optimality, because a strong feasible incumbent is what enables real decision support and strengthens the downstream exact search.

Modern branch-and-bound and branch-and-cut solvers address this difficulty by combining multiple layers of algorithmic intelligence, including presolve, cutting planes, branching strategies, node selection, and primal heuristics \cite{achterberg2009constraint,gleixner2021miplib}. Among these ingredients, primal heuristics are especially important because they determine whether a promising relaxation can be converted quickly enough into a useful incumbent \cite{berthold2013measuring}. Classical mechanisms such as the feasibility pump, local branching, and relaxation-induced neighborhood search all exploit the same underlying insight: the relaxation already contains structural information that should guide the combinatorial search \cite{fischetti2005feasibility,fischetti2003local,danna2005exploring}. Yet these heuristics also reveal a common limitation. When the feasible mixed-integer region is thin or separated from attractive fractional points by concentrated pockets of violation, local rounding and neighborhood restriction can become brittle, repeatedly revisiting similar infeasible patterns without reliably closing the final feasibility gap.

First-order optimization offers a complementary advantage on the relaxation side of the problem. Primal-dual hybrid gradient (PDHG) methods use sparse matrix-vector operations and simple projections, which makes them attractive on large and sparse relaxations for which repeated factorizations are expensive or memory intensive \cite{chambolle2011first,applegate2021practical,lu2023cupdlp}. Recent large-scale linear-programming systems show that carefully engineered primal-dual first-order methods can be implemented efficiently at industrial scale \cite{applegate2021practical,lu2023cupdlp}. However, first-order scalability does not by itself solve the mixed-integer difficulty. A PDHG trajectory may move efficiently toward a strong relaxation basin and still provide no reliable mechanism for crossing into strict integrality and linear feasibility. In other words, large-scale relaxation search and strict-feasibility recovery are related but distinct algorithmic tasks. The present paper focuses on their interface.

Our starting point is that this interface should be treated as a hybrid intelligent search problem rather than as a purely local rounding task. Hard MIP instances often require occasional non-local transitions, controlled commitment to discrete structure, and explicit feasibility repair. The intuition behind barrier-crossing moves can be motivated from several sources, including annealing theory and quantum optimization, where barrier width as well as barrier height affects the probability of escape \cite{kirkpatrick1983optimization,kadowaki1998quantum,das2008colloquium}. We do not use that literature to claim quantum hardware advantage. Instead, we use it narrowly, as an algorithmic design cue suggesting that structured non-local perturbations may help when attractive fractional basins are separated from feasible incumbents by thin but difficult violation barriers \cite{johnson2011quantum,ronnow2014defining,denchev2016computational}. For an \emph{Expert Systems with Applications} audience, the key point is therefore not the metaphor itself, but the construction of an interpretable and auditable hybrid search architecture for a practically important decision-optimization bottleneck.

Based on this perspective, we propose a hybrid population PDHG heuristic, abbreviated HP-PDHG. The method maintains a population of PDHG trajectories on the continuous relaxation, applies a barrier-aware non-local perturbation inspired by WKB-style tunneling, gradually projects integer-constrained coordinates toward the lattice, and finally subjects candidate solutions to feasibility-aware repair before they can update the incumbent. The current implementation also includes an optional feasibility-first two-phase controller that emphasizes constraint satisfaction early in the run and then returns to a more balanced search. The resulting framework is fully classical, numerically transparent, and compatible with sparse large-scale optimization. In the stored implementation and archived result files, the proposed solver still appears under the legacy label \emph{Quantum Pop-PDHG}; in this manuscript we emphasize the standard optimization interpretation and refer to it as HP-PDHG.

The paper makes four concrete contributions. First, it introduces an interpretable hybrid heuristic architecture that couples first-order relaxation search with barrier-aware perturbation, progressive integrality projection, and conservative repair. Second, it rebuilds the empirical pipeline around audited MPS parsing, original constraint senses, explicit post-repair feasibility checks, and official MIPLIB 2017 reference values, thereby removing permissive feasibility accounting inherited from earlier project iterations. Third, it demonstrates on a four-instance audited benchmark and an 80-run robustness study that HP-PDHG reliably closes the final gap between strong fractional trajectories and strict mixed-integer feasibility. Fourth, it extends the evidence to a twelve-instance official MIPLIB 2017 subset and a representative comparison with GA, DE, SHADE, and SL-PSO, then uses ablation and a 390-run sensitivity analysis to separate the truly decisive modules from the merely supportive ones.

The broader significance of the paper is methodological. It argues that feasibility recovery for MIP is a worthwhile expert-search problem in its own right: one that benefits from interpretable operator design, auditable evaluation, and careful separation between explanatory metaphor and actual mechanism. This positioning is especially important for ESWA, where contributions are expected to combine intelligent algorithm design with transparent empirical validation rather than rely on novelty in naming alone. The rest of the paper follows that principle. Section~\ref{sec:preliminaries} formalizes the feasibility gap addressed in this work. Section~\ref{sec:related_work} situates HP-PDHG within the literature on MIP heuristics, first-order optimization, learning-guided search, and barrier-crossing methods. Section~\ref{sec:theory} summarizes the analytical rationale behind population search, non-local acceptance, and progressive integrality projection. Section~\ref{sec:methodology} presents the full framework, Section~\ref{sec:experiments} reports the audited empirical evaluation, and Section~\ref{sec:conclusion} concludes with limitations and future directions.
